% !TeX program = LuaLaTeX
% Copyright (C) 2020 Roberto Giacomelli
% Barracuda manual, main TeX source file

\documentclass[11pt,a4paper]{article}
\usepackage{fontspec}
\usepackage{geometry}
\usepackage{fancyvrb}
\usepackage{graphicx}
\usepackage{booktabs}
\usepackage{array}
\usepackage{hyperref}

\newcolumntype{C}{>{\ttfamily}c}
\newcolumntype{L}{>{\ttfamily}l}

% special macro for manual typesetting









\hypersetup{
hidelinks,
linktoc = all,
pdfinfo={
    Title={the Barracuda manual},
    Subject={Barcode printing package},
    Author={Roberto Giacomelli},
    Keywords={Barcode EAN Code128 Lua}
}}

\setmainfont{Libertinus Serif}
\setmonofont[Scale=0.82]{Fira Mono}
\fvset{
    fontsize=\small,
    frame=single,
    labelposition=topline,
    framesep=5pt
}
\geometry{
    left=32mm,
    right=40mm,
    top=22mm,
    bottom=28mm
}

\author{Roberto Giacomelli\\\small email: \url{giaconet.mailbox@gmail.com}}
\title{the \code{barracuda} manual\\[1ex]
\small \url{https://github.com/robitex/barracuda}}
\date{\small Date \brcdkey{date} --- Version \brcdkey{version} --- Beta stage}

\begin{document}
\maketitle

\abstract{%
Welcome to the \brcd{} software project devoted to barcode printing.

This manual shows you how to print barcodes in your \TeX{} documents and how to
export such graphic content to an external file, using \brcd{}.

\brcd{} is written in Lua programming language and is free software released
under the GPL 2 License.%
}

\tableofcontents
\newpage


\section{Getting started}
\label{secStart}

\subsection{Introduction}
\label{secIntro}

Barcode symbols are usually a sequence of vertical lines representing encoded
data that can be retrived with special laser scanner or more simpler with a
smartphone running dedicated apps. Almost every store item has a label with a
printed barcode for automatic identification purpose.

So far, \brcd{} supported symbologies are as the following:
\begin{itemize}
\item Code 39,
\item Code 128,
\item EAN family (ISBN, ISSN, EAN 8, EAN 13, and the add-ons EAN 2 and EAN 5),
\item ITF 2of5, interleaved Two of Five ((ITF14, i2of5 general)).
\end{itemize}

The package provides different output graphic format. At the moment they are:
\begin{itemize}
\item PDF Portable Document Format (a modern \TeX{} engine is required),
\item SVG Scalable Vector Graphic.
\end{itemize}

The name \brcd{} is an assonance to the name Barcode. I started the project back
in 2016 for getting barcode in my \TeX{} generated PDF documents, studying the
Lua\TeX{} technology such as direct \emph{pdfliteral node} creation.

At the moment \brcd{} is in \emph{beta} stage. In this phase the Lua API can
change respect to the result of development research.


\subsection{Manual Content}

The manual is divided into five part. In part~\ref{secIntro} introduces the
package and gives to the user a proof of concept to how to use it. The next
parts present detailed information about option parameter of each barcode
symbology and methods description to change the \emph{module} width of a EAN-13
barcode. It's also detailed how the Lua code works internally and how to
implement a barcode symbology not already included in the package.

The plan of the manual is (but some sections are not completed yet):
\begin{description}
\item[Part 1:] Getting started
\begin{itemize}
	\item general introduction \( \to \) \pageref{secIntro}
	\item print your first barcode \( \to \) \pageref{secEnter}
	\item installing \brcd{} on your system \( \to \) \pageref{secInstall}
\end{itemize}

\item[Part 2:] \LaTeX{} packages
\begin{itemize}
	\item \brcd{} \LaTeX{} package \( \to \) \pageref{secLaTeXPkg}
\end{itemize}

\item[Part 3:] Barcode Reference and Parameters
\begin{itemize}
	\item encoder identification rule \( \to \) \pageref{secEncName}
	\item barcode symbologies reference \( \to \) \pageref{secBcRef}
\end{itemize}

\item[Part 4:] Advanced Work with \brcd{}
\begin{itemize}
	\item Lua framework description \( \to \) \pageref{secFramework}
	\item API reference \( \to \) \pageref{secAPI}
	\item \code{ga} specification \( \to \) \pageref{secGA}
\end{itemize}

\item[Part 5:] Real examples
\begin{itemize}
	\item working example and use cases \( \to \) \pageref{secExample}
\end{itemize}
\end{description}


\subsection{Required knowledge and useful resources}

\brcd{} is a Lua package that can be executed by any Lua interpreter. To use it,
it's necessary a minimal knowledge of Lua programming language and a certain
ability with the terminal of your computer system in order to run command line
task or make software installation.

It's also possible to run \brcd{} directly within a \TeX{} source file, and
compile it with a suitable typesetting engine like Lua\TeX{}. In this case a
minimal \TeX{} system knowledge is required. As an example of this workflow you
simply can look to this manual because itself is typesetted with LuaLa\TeX{},
running \brcd{} to include barcodes as a vector graphic object.

A third way is to use the \LaTeX{} package \code{barracuda.sty} with its high
level macros. A minimal knowledge of the \LaTeX{} format is obviously required.

Here is a collection of useful learning resources:
\begin{description}
\item[Lua:] to learn Lua the main reference is the book called PIL, Programming
in Lua from one of the language's Author Roberto Ierusalimschy.
\item[\LaTeX:] \dots
\item[Lua\TeX:] \dots
\end{description}


\subsection{Running Barracuda}
\label{secEnter}

The starting point to work with \brcd{} is always a plain text file with some
code processed by a command line program with a Lua interpreter.

The paradigm of \brcd{} is the Object Oriented Programming. Generally speaking
every object must be created with a function called \emph{costructor} and every
action must be run calling a \emph{method} of it.

In this section you'll take a taste of \brcd{} coding in three different
execution context: a Lua script, a Lua\TeX{} document and a \LaTeX{} source file
using the macro package \code{barracuda.sty} providing an high level interface
to Lua code.

High level package like \code{barracuda.sty} make to write Lua code unnecessary.
It will be always possible return to Lua code in order to resolve complex
barcode requirements.


\subsubsection{A Lua script}

As a practical example to produce an EAN~13 barcode, open a text editor of your
choice on an empty file and save it as \code{first-run.lua} with the content of
the following two lines of code:
\medskip
\begin{Verbatim}[label=\footnotesize\code{first-run.lua}]
local barracuda = require "barracuda"
barracuda:save("ean-13", "8006194056290", "my_barcode", "svg")
\end{Verbatim}

What you have done is to write a \emph{script}. If you have installed a Lua
interpreter along with \brcd{}, open a terminal and run it with the command:
\begin{Verbatim}
$ lua first-run.lua
\end{Verbatim}

You will see in the same directory of your script, appearing a new file called
\code{my\_barcode.svg} with the drawing:
\begin{center}
\includegraphics{image/8006194056290}
\end{center}

Coming back to the script first of all, it's necessary to load the library
\brcd{} with the standard Lua function \code{require()} that returns an
object---more precisely a reference to a table where are stored all the
package machinery.

With the second line of code, an EAN~13 barcode is saved as
\code{my\_barcode.svg} using the method \code{save()} of the \brcd{} object. The
\code{save()} method takes in order the barcode symbology identifier called
\emph{treename}, an argument as a string or as a whole number that represents
data to be encoded, the output file name and the optional output format. With a
fifth optional argument we can pass options to the barcode encoder as a Lua
table.

Each encoder has an own identifier called treename explained at
section~\ref{secEncName}. In short, in \brcd{} we can build more encoders of the
same symbology with different parameters.


\subsubsection{A Lua\TeX{} source file}

\brcd{} can also runs with Lua\TeX{} and any others Lua powered \TeX{}
engines. The source file is a bit difference respect to the previuos script: the
Lua code lives inside the argument of a \verb=\directlua= primitive, moreover we
must use an horizontal box register as output destination.
\begin{Verbatim}
% !TeX program = LuaTeX
\newbox\mybox
\directlua{
	local require "barracuda"
	barracuda:hbox("ean-13", "8006194056290", "mybox")
}\leavevmode\box\mybox
\bye
\end{Verbatim}

The method \code{hbox()} works only with Lua\TeX{}. It takes three\footnote{A
fourth argment is optional as a table with user defined barcode parameters.}
arguments: encoder \emph{treename}, encoding data as a string, the \TeX{}
horizontal box name.


\subsubsection{A Lua\LaTeX{} source file}

\LaTeX{} working minimal example would be:
\begin{Verbatim}
% !TeX program = LuaLaTeX
\documentclass{article}
\usepackage{barracuda}
\begin{document}
\barracuda{ean-13}{8006194056290}
\end{document}
\end{Verbatim}


\subsection{A more deep look}

\brcd{} is designed to be modular and flexible. For example it is possible to
draw different barcodes on the same canvas or tune barcode parameters. 

The main workflow to draw a barcode object reveals more details on internal
structure. In fact, to draw an EAN~13 barcode we must do at least the following
steps:
\begin{enumerate}
\item load the library,
\item get a reference to the \code{Barcode} abstract class,
\item build an \code{ean} encoder of the variant \code{13},
\item build an EAN~13 symbol passing data to a costructor,
\item get a reference to a new canvas object,
\item draw barcode on the canvas object,
\item get a reference of the driver object,
\item print the graphic material saving an external \code{svg} file.
\end{enumerate}

Following that step by step procedure the corresponding code is translated in
the next listing:
\begin{Verbatim}
-- lua script
local barracuda = require "barracuda" -- step 1
local barcode = barracuda:barcode() -- step 2

local ean13, err_enc = barcode:new_encoder("ean-13") -- step 3
assert(ean13, err_enc)

local symb, err_symb = ean13:from_string("8006194056290") -- step 4
assert(symb, err_symb)

local canvas = barracuda:new_canvas() -- step 5
symb:append_ga(canvas) -- step 6

local driver = barracuda:get_driver() -- step 7
local ok, err_out = driver:save("svg", canvas, "my_barcode") -- step 8
assert(ok, err_out)
\end{Verbatim}

Late the manual will give objects and methods references at
section~\ref{secAPI}.


\subsection{Installing}
\label{secInstall}

\subsubsection{Installing for Lua}

Manually copy \code{src} folder content to a suitable directory of your system
that is reachable to the system Lua interpreter.

\subsubsection{Installing for TeX Live}

If you have TeX Live installed from CTAN or from DVD TeX Collection, before any
modification to your system check if the package is already installed looking
for \emph{installed} key in the output of the command:
\begin{Verbatim}
$ tlmgr show barracuda
\end{Verbatim}

If `barracuda` is not present, run the command:
\begin{Verbatim}
$ tlmgr install barracuda
\end{Verbatim}

If you have installed TeX Live via Linux OS repository try your distribution's
package management system running a software update.

It's also possible to install the package manually:
\begin{enumerate}
\item Grab the sources from CTAN or \url{https://github.com/robitex/barracuda}.
\item Unzip it at the root of one or your TDS trees (local or personal).
\item You may need to update some filename database after this, see your \TeX{}
distribution's manual for details.
\end{enumerate}


\section{Barracuda \LaTeX{} Package}
\label{secLaTeXPkg}

The \LaTeX{} package delivered with \brcd{} is still under an early stage of
development. The only macro available is \verb=\barracuda[option]{encoder}{data}=.
A simple example is the following source file for Lua\LaTeX{}:
\begin{Verbatim}
% !TeX program = LuaLaTeX
\documentclass{article}
\usepackage{barracuda}
\begin{document}
\leavevmode
\barracuda{code39}{123ABC}\\
\barracuda{code128}{123ABC}
\end{document}
\end{Verbatim}

Every macro \verb=\barracuda= typesets a barcode symbol with the encoder defined
in the first argument, encoding data defined by the second.


\section{Barcode parameters}
\label{secBarcodeReference}

\subsection{Encoder treename}
\label{secEncName}

TODO

\subsection{Barcode Reference}
\label{secBcRef}

TODO

\section{Developer zone}

\subsection{The Barracuda Framework}
\label{secFramework}

The \brcd{} package framework consists in indipendent modules: a barcode class
hierarchy encoding a text into a barcode symbology; a geometrical library called
\code{libgeo} modelling several graphic objects; an encoding library for the
\code{ga} format (graphic assembler) and several driver to \emph{print} a
\code{ga} stream into a file or in a \TeX{} hbox register.

To implement a barcode encoder you have to write a component called
\emph{encoder} defining every parameters and implementing the encoder builder,
while a driver must understand ga opcode stream and print the corresponding
graphic object.

Every barcode encoder come with a set of parameters, some of them can be
reserved and can't be edit after the encoder was build. So, you can create many
instances of the same encoder for a single barcode type, with its own parameter
set.

The basic idea is getting faster encoders, for which the user may set up
paramenters at any level: barcode abstract class, encoder globally, down to a
single symbol object.

The Barcode class is completely indipendent from the ouput driver and viceversa.

\subsection{API reference of Lua modules}
\label{secAPI}

TODO

\subsection{\code{ga} specification}
\label{secGA}

This section defines and explains with code examples the \code{ga} instruction
stream. \code{ga} stands for \emph{graphic assembler}, a sort of essential
language that describes geometrical object like lines and rectangles mainly for
a barcode drawing library on a cartesian plane \( (O, x, y) \).

The major goal of any \brcd{} encoder is to create the \code{ga} stream
corresponding to a vector drawing of a barcode symbol.

In details, a \code{ga} stream is a numeric sequence that like a program defines
what must be draw. It is not a fully binary sequence---which is a byte stream
and ideally is what a \code{ga} stream would really be---but a sequence of
integers or floating point numbers.

In Lua this is very easy to implement. Simply append a numeric value to a table
that behave as an array. Anyway \code{ga} must be basically a binary format
almost ready to be sent or received by means of a network channel.

In the Backus–Naur form a valid \code{ga} stream grammar is described by the
following code:
\begin{Verbatim}
<valid ga stream> ::= <instructions>
<instructions> ::= <instruction>
                 | <instruction> <instructions>
<instruction> ::= <opcode>
                | <opcode> <operands>

<opcode> ::= <state>
           | <object>
           | <func>
<state> ::= 1 .. 31; graphic properties
<object> ::= 32 .. 239; graphic objects
<func> ::= 240 .. 255; functions

<operands> ::= <operand>
             | <operand> <operands>
<operand> ::= <len>
            | <coord>
            | <qty>
            | <char seq>
            | <enum>
            | <abs>
            | <points>
            | <bars>

<len> ::= f64; unit measure scaled point sp = 1/65536pt
<coord> ::= f64; unit measure scaled point sp = 1/65536pt
<qty> ::= u64
<char seq> ::= <chars> 0
<chars> ::= <char>
          | <char> <chars>
<char> ::= u64
<enum> ::= u8
<abs> ::= f64
<points> ::= <point>
           | <point> <points>
<point> ::= <x coord> <y coord>
<x coord> ::= <coord> 
<y coord> ::= <coord>
<bars> ::= <bar>
         | <bar> <bars>
<bar> := <coord> <len>

; u8 unsigned 8 bit integer
; u64 unsigned 64 bit integer
; f64 floating point 64 bit number
\end{Verbatim}

Every \code{<instruction>} changes the graphic state, for instance the current
line width, or defines a graphic object, depending on the \code{opcode} value.
Every coordinate or dimension must be expressed as \emph{scaled point}, the
basic unit of measure of \TeX{} equivalent to \( 1/65536\, \)pt.

For example, the \code{opcode} for the \code{<linewidth>} operation is 1, while
for the \code{<hline>} operation is 33. An horizontal line 6pt width from the
point (0pt, 0pt) to the point (32pt, 0pt) is represented by this \code{ga}
stream:
\begin{Verbatim}
1 393216 33 0 2097152 0
\end{Verbatim}

Introducing \code{<mnemonic code>} in \code{<opcode>} place and separate the
operations in a multiline fashion, the same sequence become more readable and
more similar to an assembler listing:
\begin{Verbatim}
<linewidth> 393216
<hline> 0 2097152 0
\end{Verbatim}

To prove and visualize the meaning of the stream, we can simply use the native
graphic driver of \brcd{} compiling this Lua\TeX{} source file:
\begin{Verbatim}
% !TeX program = LuaTeX
\newbox\mybox
\directlua{
    local barracuda = require "barracuda"
    local ga = {1, 393216, 33, 0, 2097152, 0}
    local drv = barracuda:get_driver()
    drv:ga_to_hbox(ga, "mybox")
}\leavevmode\box\mybox
\bye
\end{Verbatim}

The result is: \newbox\mybox
\directlua{
local ga = {1, 393216, 33, 0, 2097152, 0}
local drv = barracuda:get_driver()
drv:ga_to_hbox(ga, "mybox")
}\box\mybox

A more abstract way to write a \code{ga} stream consists in the \code{gacanvas}
class of the module \code{libgeo}. Every operation has been mapped to a method
named \code{encode\_<mnemonic\_opcode>()}:
\begin{Verbatim}
% !TeX program = LuaTeX
\newbox\mybox
\directlua{
    local barracuda = require "barracuda"
    local canvas = barracuda:new_canvas()
    local pt = 65536
    canvas:encode_linewidth(6*pt)
    canvas:encode_hline(0, 32*pt, 0)
    local drv = barracuda:get_driver()
    drv:ga_to_hbox(canvas, "mybox")
    tex.print("[")
    for _, n in ipairs(canvas:get_stream()) do
        tex.print(tostring(n))
    end
    tex.print("]")
} results in \box\mybox
\bye
\end{Verbatim}

\begin{figure}
\caption{A drawing example that shows how to manually set up the bounding box of
the figure, keeping automatic computation disabled for more fast stream
processing.}
\label{figSquare}
\directlua{
    local barracuda = require "barracuda"
    local canvas = barracuda:new_canvas()
    local pt = 65536
    local side = 16*pt
    local s = side/2 - 1.5*pt
    local l = side/2 - 2*pt
    local dim = 5
    canvas:encode_linewidth(1*pt)
    canvas:encode_disable_bbox()
    for c = 0, dim do
        for r = 0, dim do
            local x, y = c*side, r*side
            canvas:encode_hline(x-l, x+l, y-s)
            canvas:encode_hline(x-l, x+l, y+s)
            canvas:encode_vline(y-l, y+l, x-s)
            canvas:encode_vline(y-l, y+l, x+s)
        end
    end
    local b1 = -s - 0.5*pt
    local b2 = dim*side + s + 0.5*pt
    canvas:encode_set_bbox(b1, b1, b2, b2)
    local drv = barracuda:get_driver()
    drv:ga_to_hbox(canvas, "mybox")
}\hfill\box\mybox\hfill{}
\end{figure}

The stream is printed beside the drawing in the output PDF file. Therefore the
same \code{ga} stream can also generate a different output, for instance a SVG
file. For this purpose execute the \code{save()} method of the \code{Driver}
class (the drawing is showed in figure~\ref{figSquare}):
\begin{Verbatim}
% !TeX program = LuaTeX
\newbox\mybox
\directlua{
    local barracuda = require "barracuda"
    local canvas = barracuda:new_canvas()
    local pt = 65536
    local side = 16*pt
    local s = side/2 - 1.5*pt
    local l = side/2 - 2*pt
    local dim = 5
    canvas:encode_linewidth(1*pt)
    canvas:encode_disable_bbox()
    for c = 0, dim do
        for r = 0, dim do
            local x, y = c*side, r*side
            canvas:encode_hline(x-l, x+l, y-s)
            canvas:encode_hline(x-l, x+l, y+s)
            canvas:encode_vline(y-l, y+l, x-s)
            canvas:encode_vline(y-l, y+l, x+s)
        end
    end
    local b1 = -s - 0.5*pt
    local b2 = dim*side + s + 0.5*pt
    canvas:encode_set_bbox(b1, b1, b2, b2)
    local drv = barracuda:get_driver()
    drv:ga_to_hbox(canvas, "mybox")
    drv:save("svg", canvas, "grid")
}\leavevmode\box\mybox
\bye
\end{Verbatim}

An automatic process updates the bounding box of the figure meanwhile the stream
is read instruction after instruction. The \code{<disable\_bbox>} operation
produce a more fast execution and the figure will get the bounding box computed
until that instruction. The \code{<set\_bbox>} operation imposes a minimal
bounding box comparing to the current figure dimensions.

The initial bounding box is simply an empty figure. As a conseguence, different
strategies can be used to optimize runtime execution, such as in the previuos
code example, where bounding box is always disabled and it is set up at the last
\code{canvas} method call. More often than not, we know the bounding box of the
barcode symbol including quiet zones.

Every encoding method of \code{gaCanvas} class gives two output result: a
boolean value called \code{ok} plus an error \code{err}. If \code{ok} is
\code{true} then \code{err} is \code{nil} and, viceversa, when \code{ok} is
\code{false} then \code{err} is a string describing the error.

The error management is a responsability of the caller. For instance, if we
decide to stop the execution this format is perfectly suitable for the Lua
function \code{assert()}, otherwise we can explicity check the output pair:
\begin{Verbatim}
local pt = 65536
assert(canvas:encode_linewidth(6*pt)) --> true, nil
local ok, err = canvas:encode_hline(nil, 32*pt, 0)
-- ok = false
-- err = "[ArgErr] 'x1' number expected"
\end{Verbatim}


\subsubsection{\code{ga} reference}
\label{secGAtabref}

\noindent
\begin{tabular}{CLlL}
\toprule
\multicolumn{4}{l}{\textbf{Properties of the graphic state}}\\
OpCode & Mnemonic key & Graphic property & Operands\\
\midrule
 1 & linewidth & Line width      & w <len>\\
 2 & linecap   & Line cap style  & e <enum>\\
   & & & 0: Butt cap\\
   & & & 1: Round cap\\
   & & & 2: Projecting square cap\\
 3 & linejoin  & Line join style & e <enum>\\
   & & & 0: Miter join\\
   & & & 1: Round join\\
   & & & 2: Bevel join\\
\midrule
29 & enable\_bbox & Compute bounding box & -\\
30 & disable\_bbox & Do not compute bounding box & -\\
31 & set\_bbox & Overlap a bounding box & x1 y1 <point> x2 y2 <point>\\
\bottomrule
\end{tabular}

\bigskip
\noindent
\begin{tabular}{CLlL}
\toprule
\multicolumn{4}{l}{\textbf{Lines}}\\
OpCode & Mnemonic key & Graphic object & Operands\\
\midrule
32 & line  & Line            & x1 y1 <point> x2 y2 <point>\\
33 & hline & Horizontal line & x1 x2 <point> y <coord>\\
34 & vline & Vertical line   & y1 y2 <point> x <coord>\\
\bottomrule
\end{tabular}

\bigskip
\noindent
\begin{tabular}{CLlL}
\toprule
\multicolumn{4}{l}{\textbf{Group of bars}}\\
OpCode & Mnemonic key & Graphic object & Operands\\
\midrule
36 & vbar & Vertical bars & y1 <coord> y2 <coord> b <qty> xi wi <bars>\\
   & & & y1: bottom y-coord\\
   & & & y2: top y-coord\\
   & & & b: number of bars\\
   & & & xi: axis x-coord of bars number i\\
   & & & wi: width of bars number i\\
37 & hbar & Horizontal bars & x1 <coord> x2 <coord> b <qty> yi wi <bars>\\
   & & & unimplemented\\
\midrule
38 & polyline & Opened polyline & n <qty> xi yi <points>\\
   & & & n: number of points\\
   & & & xi: x-coord of point i\\
   & & & yi: y-coord of point i\\
39 & c\_polyline & Closed polyline & n <qty> xi yi <points>\\
   & & & unimplemented\\
\bottomrule
\end{tabular}

\bigskip
\noindent
\begin{tabular}{CLlL}
\toprule
\multicolumn{4}{l}{\textbf{Rectangles}}\\
OpCode & Mnemonic key & Graphic object & Operands\\
\midrule
48 & rect & Rectangle & x1 y1 <point> x2 y2 <point>\\
49 & f\_rect & Filled rectangle & x1 y1 <point> x2 y2 <point>\\
 &&& unimplemented\\
50 & rect\_size & Rectangle & x1 y1 <point> w <len> h <len>\\
 &&& unimplemented\\
51 & f\_rect\_size & Filled rectangle & x1 y1 <point> w <len> h <len>\\
 &&& unimplemented\\
\bottomrule
\end{tabular}

\bigskip
\noindent
\begin{tabular}{CLl}
\toprule
\multicolumn{3}{l}{\textbf{Text}}\\
OpCode & Mnemonic key & Graphic object/Operands\\
\midrule
130 & text & A text with several glyphs\\
    & & \ttfamily ax <abs> ay <abs> xpos ypos <point> c <chars>\\
\midrule
131 & text\_xspaced & A text with glyphs equally spaced on its vertical axis\\
    & & \ttfamily x1 <coord> xgap <len> ay <abs> ypos <coord> c <chars>\\
\midrule
132 & text\_xwidth & Glyphs equally spaced on vertical axis between two x coordinates\\
    & & \ttfamily ay <abs> x1 <coord> x2 <coord> y <coord> c <chars>\\
\midrule
140 & \_text\_group & Texts on the same baseline \\
& & \ttfamily ay <abs> y <coord> n <qty> xi <coord> ai <abs> ci <chars> \\
& & unimplemented\\
\bottomrule
\end{tabular}

\subsection{\code{Vbar} class}

This section show you how to draw a group of vertical lines, the main component
of every 1D barcode symbol. In the \brcd{} jargon a group of vertical lines is
called \code{Vbar} and is defined by a flat array of pair numbers: the first one
is the x-coordinate of the bar while the second is the its width.
Our \code{Vbar} is represented by this array:
\begin{Verbatim}
-- {   x1,    w1,      x2,    w2,     x3,     w3}
   {1*mod, 2*mod, 3.5*mod, 1*mod, 5.5*mod, 1*mod}
\end{Verbatim}
We have three bars based on the fixed lenght \code{mod} represented as:
\begin{center}

\end{center}





In the module \code{libgeo} there is a class called \code{Vbar}

\subsection{\code{ga} code examples}

To support a better understanding and provide some \code{ga} stream examples,
this section discusses some code examples, each of which must be compiled with
Lua\TeX{}.

\subsubsection{Example number 1}

Suppose we want to draw a simple rectangle. In the \code{ga} reference at
section~\ref{secGAtabref} there is a dedicated instruction \code{<rect>}.
Let's give it a try:
\begin{Verbatim}
% !TeX program = LuaTeX
\newbox\mybox
\directlua{
    local barracuda = require "barracuda"
    local pt = 65536
    local side = 36*pt
    local ga = {48, 0, 0, 2*side, side}
    local drv = barracuda:get_driver()
    drv:ga_to_hbox(ga, "mybox")
}\leavevmode\box\mybox
\bye
\end{Verbatim}

Dealing with low level \code{ga} stream is not necessary. We can use a
\code{gaCanvas} object:
\begin{Verbatim}
% !TeX program = LuaTeX
\newbox\mybox
\directlua{
    local barracuda = require "barracuda"
    local pt = 65536
    local side = 36*pt
    local canvas = barracuda:new_canvas()
    canvas:encode_rect(0, 0, 2*side, side)
    local drv = barracuda:get_driver()
    drv:ga_to_hbox(canvas, "mybox")
}\leavevmode\box\mybox
\bye
\end{Verbatim}

In both cases the drawing looks like: \directlua{
    local barracuda = require "barracuda"
    local pt = 65536
    local side = 36*pt
    local canvas = barracuda:new_canvas()
    canvas:encode_rect(0, 0, 2*side, side)
    local drv = barracuda:get_driver()
    drv:ga_to_hbox(canvas, "mybox")
}\box\mybox

\subsubsection{Example number 2}

A more complex drawing is a chess. Let's begin to draw a single cell with a
square 1cm width:
\begin{Verbatim}
% !TeX program = LuaTeX
\newbox\mybox
\directlua{
    local barracuda = require "barracuda"
    local drv = barracuda:get_driver()
    local mm = drv.mm
    local s, t = 10*mm, 2*mm
    local canvas = barracuda:new_canvas()
    canvas:encode_linewidth(t)
    canvas:encode_rect(t/2, t/2, s-t/2, s-t/2)
    drv:ga_to_hbox(canvas, "mybox")
}\leavevmode\box\mybox
\bye
\end{Verbatim}

and repeat the game for the chess grid:
\begin{Verbatim}
% !TeX program = LuaTeX
\newbox\mybox
\directlua{
    local barracuda = require "barracuda"
    local drv = barracuda:get_driver()
    local mm = drv.mm
    local s, t = 10*mm, 2*mm
    local canvas = barracuda:new_canvas()
    canvas:encode_linewidth(t)
    for row = 1, 5 do
        for col = 1, 5 do
            local l = (row + col)/2
            if l == math.floor(l) then
                local x = (col - 1) * s
                local y = (row - 1) *s
                canvas:encode_rect(x+t/2, y+t/2, x+s-t/2, y+s-t/2)
            end
        end
    end
    drv:ga_to_hbox(canvas, "mybox")
}\leavevmode\box\mybox
\bye
\end{Verbatim}

The final drawing is showed here: \directlua{
    local barracuda = require "barracuda"
    local drv = barracuda:get_driver()
    local mm = drv.mm
    local s, t = 10*mm, 2*mm
    local canvas = barracuda:new_canvas()
    canvas:encode_linewidth(t)
    for row = 1, 5 do
        for col = 1, 5 do
            local l = (row + col)/2
            if l == math.floor(l) then
                local x = (col - 1) * s
                local y = (row - 1) *s
                canvas:encode_rect(x+t/2, y+t/2, x+s-t/2, y+s-t/2)
            end
        end
    end
    drv:ga_to_hbox(canvas, "mybox")
}\box\mybox


\subsubsection{Example number 3}

A drawing of a zig zag staircase can be represented by a \code{ga} stream with
a \code{<polyline>} operation. The \code{gaCanvas} method we have to call is
\code{encode\_polyline()} that accept a Lua table as a flat structure with the
coordinates of every point of the polyline:
\begin{Verbatim}
{ x1, y1, x2, y2, ..., xn, yn }
\end{Verbatim}

It is what we do with this code:
\begin{Verbatim}
% !TeX program = LuaTeX
\newbox\mybox
\directlua{
    local barracuda = require "barracuda"
    local pt = 65536
    local side = 18*pt
    local dim = 5
    local x, y = 0, 0
    local point = {x, y}
    local i = 3
    for _ = 1, dim do
        y = y + side
        point[i] = x; i = i + 1
        point[i] = y; i = i + 1
        x = x + side
        point[i] = x; i = i + 1
        point[i] = y; i = i + 1
    end
    local canvas = barracuda:new_canvas()
    canvas:encode_linewidth(2.5*pt)
    canvas:encode_polyline(point)
    local drv = barracuda:get_driver()
    drv:ga_to_hbox(canvas, "mybox")
}\leavevmode\box\mybox
\bye
\end{Verbatim}

The result is: \directlua{
    local pt = 65536
    local side = 18*pt
    local dim = 5
    local x, y = 0, 0
    local point = {x, y}
    local i = 3
    for _ = 1, dim do
        y = y + side
        point[i] = x; i = i + 1
        point[i] = y; i = i + 1
        x = x + side
        point[i] = x; i = i + 1
        point[i] = y; i = i + 1
    end
    local canvas = barracuda:new_canvas()
    canvas:encode_linewidth(2.5*pt)
    canvas:encode_polyline(point)
    local drv = barracuda:get_driver()
    drv:ga_to_hbox(canvas, "mybox")
}\leavevmode\box\mybox

A feature of \code{encode\_<opcode>} methods is a \emph{polymorfic} behaviour.
They accept different types for own arguments. \code{encode\_polyline} is not an
exception: it accept also a \code{Polyline} object of the \code{libgeo} module
instead of a flat array of coordinates. This feature is showed by the following
code that re-implement the previuos version:
\begin{Verbatim}
% !TeX program = LuaTeX
\newbox\mybox
\directlua{
    local barracuda = require "barracuda"
    local pt = 65536
    local side = 18*pt
    local dim = 5
    local Polyline = barracuda:libgeo().Polyline
    local pl = Polyline:new()
    pl:add_point(0, 0)
    for _ = 1, dim do
        pl:add_relpoint(0, side)
        pl:add_relpoint(side, 0)
    end
    local canvas = barracuda:new_canvas()
    canvas:encode_linewidth(2.5*pt)
    canvas:encode_polyline(pl)
    local drv = barracuda:get_driver()
    drv:ga_to_hbox(canvas, "mybox")
}\leavevmode\box\mybox
\bye
\end{Verbatim}

Pretty sure that this new version is more clear and intuitive.


%\subsubsection{Example number 4}

%A polyline that represents a path of ... Hilbert


%\subsubsection{Example number 5}

%Text pyramid


\section{Example and use cases}
\label{secExample}

TODO

\end{document}
